\documentclass[11pt]{scrartcl}
\usepackage[a4paper,margin=2.6cm]{geometry}
\usepackage{../common/rlstudy}
\title{Extended reading --- Week 10: SAC, TD3, CrossQ}
\date{\today}
\begin{document}
\maketitle
\section*{Read with a question in mind}
Start with the week's teaching guide and its small worked example. Then use
these readings to connect the algorithm to its assumptions and evidence.
The explanations and investigations below are course commentary, not
quotations. Read the source before reusing a theorem or a reported result.
For each paper, record the objective, observation model, claimed result,
conditions and one condition your local experiment does not reproduce.

\section*{1. Lillicrap et al. (2015/2016): Continuous Control with Deep Reinforcement Learning}
\textbf{Primary source or book record:} \url{https://arxiv.org/abs/1509.02971}.

\paragraph{What to learn.} DDPG combines a deterministic actor-critic with replay and slowly changing targets for continuous actions.

\paragraph{What the result requires.} It is not the first possible deep continuous-control method, and its components have prior histories. Replay-distribution actor updates require care when comparing them with an exact start-state objective.

\paragraph{Notebook investigation.} Trace the actor derivative through Q with respect to action, while keeping critic parameters fixed for that actor update. Distinguish this derivative from the critic regression gradient.

\section*{2. Fujimoto, van Hoof and Meger (2018): Addressing Function Approximation Error in Actor-Critic Methods}
\textbf{Primary source or book record:} \url{https://arxiv.org/abs/1802.09477}.

\paragraph{What to learn.} TD3 combines clipped double critics, delayed actor updates and target-policy smoothing.

\paragraph{What the result requires.} The minimum can introduce pessimism; it is not universally a low-variance unbiased estimate. Published gains depend on task and protocol, not a twofold improvement on every benchmark.

\paragraph{Notebook investigation.} Use \texttt{td3\_target}. Apply noise clipping, action-bound clipping and a terminal mask in that order before comparing the two critic values.

\section*{3. Haarnoja et al. (2018): Soft Actor-Critic and Algorithms and Applications}
\textbf{Primary source or book record:} \url{https://arxiv.org/abs/1812.05905}.

\paragraph{What to learn.} SAC optimises an entropy-regularised objective. The later formulation adapts temperature through a target-entropy constraint; the original paper is at arXiv:1801.01290.

\paragraph{What the result requires.} Automatic temperature leaves a target entropy and optimisation choices. Entropy depends on action coordinates and, for a continuous distribution, need not be positive. The tanh Jacobian is part of the action density.

\paragraph{Notebook investigation.} Use \texttt{tanh\_gaussian\_log\_prob} and \texttt{auto\_temperature\_loss}. Keep sampled actions fixed when checking the intended direction of a temperature update.

\section*{4. Bhatt et al. (2024): CrossQ}
\textbf{Primary source or book record:} \url{https://arxiv.org/abs/1902.05605}.

\paragraph{What to learn.} Examines batch normalisation and joint processing of current and successor critic inputs, achieving strong sample efficiency without target networks in the studied setting.

\paragraph{What the result requires.} The mechanism is more specific than adding batch norm after every linear layer. Successor targets still require the intended gradient stopping. Neither a fixed speedup nor a universal replacement guarantee follows.

\paragraph{Notebook investigation.} Inspect the primary implementation before designing an ablation: which samples share normalisation statistics, and which gradients stop? The local \texttt{polyak\_update} exercise teaches a different target-network mechanism.

\section*{5. Lee et al. (2025): SimBa: Simplicity Bias for Scaling Up Parameters in Deep Reinforcement Learning}
\textbf{Primary source or book record:} \url{https://arxiv.org/abs/2410.09754}.

\paragraph{What to learn.} Studies an architecture using observation normalisation, residual feed-forward blocks and layer normalisation to improve parameter scaling.

\paragraph{What the result requires.} Do not mix this original architecture with later variants involving additional normalisations, resets or distributional losses. Its published benchmark evidence is not a universal current-best claim.

\paragraph{Notebook investigation.} Draw the original residual block from the paper. Identify the normalisation axis and where observations are standardised; then describe a controlled architecture-only comparison.

\section*{Synthesis: evidence you should be able to show}
\begin{enumerate}
\item Reproduce one numerical calculation before running a long experiment.
\item Name the precise objective, sampling distribution and terminal convention.
\item Compare the local implementation with one source: list a shared mechanism
and one simplification that prevents calling it a full reproduction.
\item Separate an algebraic identity, an assumption-dependent theorem and an
empirical result. A plot alone does not prove a theorem.
\item Save the notebook with seed, parameters, axes, denominators and a short
explanation of an unexpected result. Keep evaluation data separate from tuning.
\end{enumerate}
Return to the teaching guide's notebook investigation for the local function
names, expected checks and interpretation questions. Source code and notebooks
are available together in the course repository; the web reader is the entry
point and the downloaded notebook is the execution workspace.
\section*{Statistics reading bridge}
\textbf{The minimum of two unbiased critics can be biased.} Read bias and expectation of transformed random variables, then enumerate the twin-error fixture. Return to TD3 and SAC primary sources for why their particular targets use the minimum; do not call that minimum an unbiased value estimate.

\textbf{Checked textbook locations:} Siegrist \S~\S~4.1, 7.1, PDF pp.439, 896. \url{https://stats.libretexts.org/@go/page/10189}. These locations support the probability and estimation concepts; the algorithm-specific results remain attributed to the primary readings above. The course source catalog distinguishes the original textbook from its supplied instructor manual and records PDF-viewer pagination.

\end{document}
